% !TeX root = ../main.tex
% !TeX encoding = UTF-8

\chapter*{Appendix B: The Digital Alibi Assessment Protocol, Board Checklist}
\addcontentsline{toc}{chapter}{Appendix B: The Digital Alibi Assessment Protocol, Board Checklist}
\markboth{Appendix B: The Digital Alibi Assessment Protocol, Board Checklist}{Appendix B: The Digital Alibi Assessment Protocol, Board Checklist}

This checklist is a preparatory document for the audit committee or board risk committee ahead of commissioning a Digital Alibi Assessment. It does not replace the assessment. It establishes whether the board is ready to commission one and what the committee should require of it.

\textbf{1. AI systems inventory confirmation.} Does a current inventory of AI systems exist, including vendor-embedded and third-party AI? When was it last refreshed, and by whom? Has a systematic shadow AI surfacing exercise been conducted, or does the inventory rest on self-declaration by business units?

\textbf{2. Material decision identification.} Which decision categories (credit, underwriting, claims, hiring, procurement, compliance determinations) involve AI systems materially? Can the organisation identify, for a named recent period, the specific material decisions each system produced or influenced? Has the board defined the materiality threshold, or does it exist only by implication?

\textbf{3. Four-element gap assessment.} For a sample of recent material AI-assisted decisions, can the organisation produce, from contemporaneous records: Decision Identity, which system, at which version, produced which output in which context? The Information Picture, what the system was given, what it retrieved, and what it was structurally unable to see? The Human Oversight Record, who reviewed the output, what they were shown, what authority they held, and what they decided? The Accountability Chain, the named individuals accountable at the moment of the decision, assigned at the time?

\textbf{4. Accountability chain verification.} Is accountability assigned to ongoing decision validity, or only to system approvals? Who owns the obligation to reassess an approved system when external conditions change? Does the chain reach the board, and is it current within the last review cycle?

\textbf{5. Retrievability.} Has retrieval of a complete decision record ever been tested against a realistic examination timeframe? Who would conduct that retrieval, and from which systems?

\textbf{6. Commissioning requirements.} The assessment the committee commissions should be independent of the function under examination, examine specific decisions rather than framework documentation, produce a confirmed governance address rather than a self-assessed one, and report its findings to the committee without internal filtering.

A committee that cannot complete sections 1 to 5 from existing records has, in that inability, its first finding. The assessment does not create the gaps it names. It converts unexamined exposure into a deliberate governance response.
